%!TeX program = lualatex
\renewcommand{\thelektion}{L02 \textperiodcentered{} Basisgrössen \& Stellenwertdarstellung}

\begin{frame}\lessontitle{Lektion 02}{Basisgrössen \& Stellenwertdarstellung}{Von römischen Münzen zum Dezimal- und Binärsystem \textperiodcentered{} Skript, Kapitel 2}\end{frame}

\begin{frame}\FDUtitle{Rückblick}
Letztes Mal: die \tib{unäre} Zahlendarstellung.
\vspace{3mm}
\begin{center}
||||||||||||||||||||||||||||||||||||||
\end{center}
\begin{itemize}\setlength\itemsep{5pt}
  \item Ein einziges Symbol, Länge wächst linear
  \item Heute: Kulturen der Antike verwendeten bereits \tib{mehrere} Symbole mit unterschiedlichem Wert
\end{itemize}
\end{frame}

{\setbeamertemplate{footline}{\darkfootline}
\begin{frame}\sectionframe[FDUgreen]{Römische Zahlen}{Sieben Basisgrössen: $I, V, X, L, C, D, M$}
\end{frame}}

\begin{frame}\FDUtitle{Basisgrössen: Römische Zahlen}
\begin{center}
Welche Zahl steckt dahinter?
\end{center}
\begin{equation*}
	CCCCCLIIII
\end{equation*}
\only<2->{
\begin{center}
554 $\rightarrow$ Gibt es eine kürzere Darstellung?
\end{center}}
\only<3->{
\begin{equation*}
	DLIV
\end{equation*}}
\only<4->{\begin{center}\color{FDUteal}\bfseries Basisgrössen lassen sich wie \textbf{Münzen} verstehen.\end{center}}
\end{frame}

\begin{frame}\FDUtitle{Die Münzen-Analogie}
Eine Sammlung von Münzen heisst \tib{Zahlendarstellung} einer Zahl $m$, wenn:
\begin{itemize}\setlength\itemsep{6pt}
  \item \tib{Korrektheit}: die Summe der Münzenwerte gleich $m$ ist,
  \item \tib{Minimalität}: keine Sammlung mit weniger Münzen existiert, deren Summe $m$ ist.
\end{itemize}
\vspace{3mm}
\begin{center}
Beispiel: $VII$ ist die minimale römische Münzendarstellung von \textit{sieben} (3 statt 7 Münzen).
\end{center}
\end{frame}

\begin{frame}\FDUtitle{Nicht immer eindeutig}
Wir haben folgende Münzen: \coin{1}\coin{2}\coin{3}\coin{5}\coin{10}
\vspace{3mm}
\begin{center}
Wie stellt man die Zahl 4 dar?
\end{center}
\only<2->{
\begin{itemize}
\item \coin{1}\coin{3}
\item \coin{2}\coin{2}
\end{itemize}}
\only<3->{
\begin{center}\color{FDUteal}\bfseries $\rightarrow$ Die Darstellung ist minimal, aber nicht eindeutig.\end{center}
}
\end{frame}

\begin{frame}{Auftrag}{Kapitel 2: Ursprünge \& Basisgrössen}
\aufgabebox{\textbf{Aufgabe 2.1}: Römisches Münzsystem um zwei Münztypen erweitern}
\vspace{3mm}
\aufgabebox{\textbf{Aufgabe 2.2}: Zeichenanzahl für 1 Milliarde im römischen System bestimmen}
\end{frame}

{\setbeamertemplate{footline}{\darkfootline}
\begin{frame}\sectionframe[FDUblue]{Stellenwertdarstellung}{Vom Dezimal- zum Binärsystem}
\end{frame}}

\begin{frame}\FDUtitle{Warum unterschiedliche Zahlensysteme?}
\begin{itemize}\setlength\itemsep{6pt}
  \item Zahlensysteme sind \enquote{willkürlich} und historisch entstanden
  \item Computer arbeiten mit binären Zahlen $\lrc{0,1}$
  \begin{itemize}
    \item Schnelleres, zuverlässigeres Rechnen
    \item \emoji{high-voltage} 1 = Strom fliesst, 0 = kein Strom
  \end{itemize}
\end{itemize}
\end{frame}

\begin{frame}\FDUtitle{Stellenwertdarstellung: Beobachtungen}
\begin{itemize}\setlength\itemsep{8pt}
    \item Alphabet des \enquote{normalen} Zahlensystems:
    $$\lrc{0,1,2,3,4,5,6,7,8,9}$$
    \begin{itemize}
    \item \textcolor{FDUblue}{10-adisches} Zahlensystem, da 10 Zeichen
    \item \textcolor{FDUteal}{Dezimalsystem} (lat.\ \textit{decem} $\rightarrow$ \enquote{zehn})
    \end{itemize}
    \item Alphabet des binären Zahlensystems:
    $$\lrc{0,1}$$
    \begin{itemize}
    \item \textcolor{FDUblue}{2-adisches} Zahlensystem, da 2 Zeichen
    \item \textcolor{FDUteal}{Binärsystem} (lat.\ \textit{binarius} $\rightarrow$ \enquote{zweifach})
    \end{itemize}
\end{itemize}
\end{frame}

\begin{frame}\FDUtitle{Dezimal $\rightarrow$ Binär: erste Beispiele}
\only<1>{
\begin{table}[H]
    \centering
    \begin{tabular}{c|c}
         \textbf{Dezimal} & \textbf{Binär}  \\\midrule\midrule
         $0$&$?_{2}$\\\hline
         $1$&$?_{2}$\\\hline
         $2$&$?_{2}$\\\hline
         $3$&$?_{2}$\\\hline
         $4$&$?_{2}$\\\hline
         ...&...
    \end{tabular}
\end{table}
}
\only<2->{
\begin{table}[H]
    \centering
    \begin{tabular}{c|c}
         \textbf{Dezimal} & \textbf{Binär}  \\\midrule\midrule
         $0$&$0_{2}$\\\hline
         $1$&$1_{2}$\\\hline
         $2$&$10_{2}$\\\hline
         $3$&$11_{2}$\\\hline
         $4$&$100_{2}$\\\hline
         ...&...
    \end{tabular}
\end{table}
}
\only<3->{\begin{center}Wie rechnet man das systematisch um? $\rightarrow$ \tib{$b$-adische Zahlensysteme}\end{center}}
\end{frame}

\begin{frame}\FDUtitle{Basisgrössen im Stellenwertsystem}
Das \textcolor{FDUteal}{Dezimalsystem} hat die Basisgrössen
$$10^0=1,\quad 10^1=10,\quad 10^2=100,\quad 10^3=1000,\ldots$$
\rule{\textwidth}{0.4pt}
Beispiel: Das Wort \textcolor{FDUteal}{$4605$} wird interpretiert als
\only<1>{
\begin{align*}
    \textcolor{FDUteal}{\underbrace{\ldots}_{?}} \underbrace{10^3}_{?} +
    \textcolor{FDUteal}{\underbrace{\ldots}_{?}} \underbrace{10^2}_{?} +
    \textcolor{FDUteal}{\underbrace{\ldots}_{?}} \underbrace{10^1}_{?} +
    \textcolor{FDUteal}{\underbrace{\ldots}_{?}} \underbrace{10^0}_{?}.
\end{align*}}
\only<2->{
\begin{align*}
    \textcolor{FDUteal}{4}\cdot \underbrace{10^3}_{1000} + \textcolor{FDUteal}{6}\cdot \underbrace{10^2}_{100} + \textcolor{FDUteal}{0}\cdot \underbrace{10^1}_{10} + \textcolor{FDUteal}{5}\cdot \underbrace{10^0}_{1}
\end{align*}
\only<3->{= vier Tausender $+$ sechs Hunderter $+$ null Zehner $+$ fünf Einer.}
}
\end{frame}

\begin{frame}\FDUtitle{Definition: $b$-adische Zahlendarstellung}
Für eine Basis $b\geq 2$ definieren wir
\begin{align*}
	(a_{n}a_{n-1}\ldots a_{1}a_{0})_{b} := a_n\cdot b^n + a_{n-1}\cdot b^{n-1} + \ldots + a_1\cdot b + a_0,
\end{align*}
wobei $a_0,\ldots,a_n \in \lrc{0,1,\ldots,b-1}$ \tib{Ziffern} genannt werden.
\vspace{3mm}
\begin{center}
Beispiel ($b=3$): $\quad 211_3 = 2\cdot 3^2 + 1\cdot 3^1 + 1\cdot 3^0 = 22$
\end{center}
\end{frame}

\begin{frame}\FDUtitle{Wichtige Beobachtung}
\begin{center}
Warum sind im $b$-adischen System nur Ziffern $0,\ldots,b-1$ erlaubt?
\end{center}
\only<2->{
\begin{itemize}\setlength\itemsep{6pt}
\item Beispiel: 13 Objekte im Dezimalsystem mit möglichst wenigen \enquote{Münzen} darstellen.
\item 10 Einer lassen sich stets durch 1 Zehner ersetzen $\rightarrow$ 9 Münzen gespart.
\begin{itemize}
    \item Münzen vorher: \coin{1}\coin{1}\coin{1}\coin{1}\coin{1}\coin{1}\coin{1}\coin{1}\coin{1}\coin{1}\coin{1}\coin{1}\coin{1}
    \item Münzen nachher: \coin{10}\coin{1}\coin{1}\coin{1}
\end{itemize}
\item Allgemein: $b$ Münzen vom Typ $b^k$ lassen sich immer durch eine Münze vom Typ $b^{k+1}$ ersetzen.
\end{itemize}}
\end{frame}

\begin{frame}\FDUtitle{Warum ist das für die Informatik wichtig?}
Die Bedingung $a_k \in \lrc{0,1,\ldots,b-1}$ bringt zwei entscheidende Vorteile:
\vspace{3mm}
\begin{enumerate}\setlength\itemsep{8pt}
  \item \textbf{Minimale Darstellung (Speichereffizienz):}\\
  Die Darstellung verbraucht so wenige Stellen (Bits) wie möglich.
  \item \textbf{Eindeutigkeit (Kanonische Form):}\\
  Jede Zahl besitzt exakt \emph{eine} Bitfolge $\rightarrow$ Computer können Werte ohne Mehrdeutigkeit vergleichen und speichern.
\end{enumerate}
\vspace{3mm}
\only<2->{
\begin{center}
  \color{FDUteal}\bfseries Das Bündeln von $b$ Einheiten zur nächsten Stelle entspricht physikalisch dem \textbf{Übertrag} (\textit{Carry-Bit}) in Prozessoren!
\end{center}
}
\end{frame}

\begin{frame}{Auftrag}{Kapitel 2: Stellenwertdarstellung}
\aufgabebox{\textbf{Aufgabe 2.3 \& 2.4}: Binärzahl $11011_2 \rightarrow$ Dezimal und Dezimal $14 \rightarrow$ Binär}
\vspace{3mm}
\aufgabebox{\textbf{Aufgabe 2.5}: Darstellung der Zahl 29 im Binärsystem mit $\{0,1\}$}
\vspace{3mm}
\challengebox{\textbf{Aufgabe 2.6}: Warum ist Basis $b=1$ nicht sinnvoll?}
\end{frame}

